% Purpose: Master Beamer presentation on AI in Education,
% Epistemology, and Conatus ad Novandum.
% Status: Complete Draft; ready for lecture/seminar delivery and review.
% Source-of-truth: doc/EduAI/abstract/main.tex
\documentclass[aspectratio=169,11pt]{beamer}

\usepackage[spanish,es-nodecimaldot]{babel}
\usepackage[utf8]{inputenc}
\usepackage[T1]{fontenc}
\usepackage{lmodern}
\usepackage{amsmath,amssymb}
\usepackage{booktabs}
\usepackage{microtype}
\usepackage{tikz}
\usetikzlibrary{arrows.meta,positioning,shapes.geometric,calc,fit,decorations.pathreplacing}

% -----------------------------------------------------------------------------
% Color Palette & Visual Theme Configuration (Modern Minimalist 16:9)
% -----------------------------------------------------------------------------
\definecolor{primaryDark}{HTML}{1A2B4C}      % Deep Navy
\definecolor{secondarySlate}{HTML}{2E4057}   % Slate Blue
\definecolor{accentGold}{HTML}{C59B27}       % Warm Gold / Ochre
\definecolor{accentTeal}{HTML}{0E7C86}       % Deep Teal
\definecolor{bgSoft}{HTML}{F8F9FA}           % Off-white soft background
\definecolor{cardBg}{HTML}{FFFFFF}           % Clean white for cards
\definecolor{textDark}{HTML}{1E252B}         % Primary reading text
\definecolor{textMuted}{HTML}{5A6570}        % Muted caption/subtitle text
\definecolor{borderSubtle}{HTML}{E2E8F0}     % Subtle card boundary

\setbeamercolor{background canvas}{bg=bgSoft}
\setbeamercolor{normal text}{fg=textDark}
\setbeamercolor{frametitle}{fg=primaryDark,bg=bgSoft}
\setbeamercolor{title}{fg=primaryDark}
\setbeamercolor{subtitle}{fg=secondarySlate}
\setbeamercolor{author}{fg=textDark}
\setbeamercolor{date}{fg=textMuted}
\setbeamercolor{item}{fg=accentTeal}
\setbeamercolor{subitem}{fg=accentGold}

\setbeamertemplate{navigation symbols}{}
\setbeamertemplate{itemize
item}{\raisebox{0.1ex}{\footnotesize$\blacktriangleright$}}
\setbeamertemplate{itemize subitem}{\raisebox{0.15ex}{\tiny$\blacksquare$}}

% Frametitle layout
\setbeamertemplate{frametitle}{
  \vspace{0.25cm}
  \begin{beamercolorbox}[wd=\paperwidth,leftskip=0.8cm,rightskip=0.8cm]{frametitle}
    {\usebeamerfont{frametitle}\textbf{\color{primaryDark}\insertframetitle}}\par
    \ifx\insertframesubtitle\@empty\else
    \vspace{0.05cm}
    {\usebeamerfont{framesubtitle}\color{accentTeal}\insertframesubtitle}\par
    \fi
  \end{beamercolorbox}
  \vspace{-0.15cm}
}

% Footline layout
\setbeamertemplate{footline}{
  \leavevmode%
  \hbox{%
    \begin{beamercolorbox}[wd=0.7\paperwidth,ht=2.2ex,dp=1.0ex,leftskip=0.8cm]{
        author in head/foot
      }
      \usebeamerfont{footline}\color{textMuted}\textbf{
        Educación en la Era de la IA
      } \textbar\ Ken Tanaka Hernandez
    \end{beamercolorbox}%
    \begin{beamercolorbox}[wd=0.3\paperwidth,ht=2.2ex,dp=1.0ex,rightskip=0.8cm]{
        date in head/foot
      }
      \hfill\usebeamerfont{footline}\color{textMuted}\insertframenumber\ /
      \inserttotalframenumber
    \end{beamercolorbox}%
  }%
  \vspace{0.15cm}
}

% Semantic Card Block
\newcommand{\conceptbox}[3]{%
  \begin{tikzpicture}
    \node[
      fill=#1!6,
      draw=#1!40,
      line width=0.7pt,
      rounded corners=3pt,
      inner sep=5.5pt,
      text width=\dimexpr\linewidth-13pt\relax,
      align=left
    ] {
      {\bfseries\footnotesize\color{#1}#2}\\[0.2em]
      \scriptsize\color{textDark}#3
    };
  \end{tikzpicture}%
}

% -----------------------------------------------------------------------------
% Presentation Metadata
% -----------------------------------------------------------------------------
\title{\huge\textbf{Educación en la era de la
IA:}\\[0.2em]\LARGE¿Cultivando a los futuros pensadores?}
\author{\textbf{Ken Tanaka Hernandez}}
\institute{Universit\`{a} degli Studi di Trieste \textbar\ OGS}
\date{\today}

% -----------------------------------------------------------------------------
% Document Body
% -----------------------------------------------------------------------------
\begin{document}

% -----------------------------------------------------------------------------
% SLIDE 1: Portada
% -----------------------------------------------------------------------------
\begin{frame}[plain]
  \vspace{1.0cm}
  \begin{beamercolorbox}[sep=8pt,left,wd=\textwidth]{title}
    \usebeamerfont{title}\inserttitle\par
  \end{beamercolorbox}
  \vspace{0.6cm}
  \begin{beamercolorbox}[sep=8pt,left,wd=\textwidth]{author}
    \usebeamerfont{author}\insertauthor\par
    \vspace{0.15cm}
    \small\color{textMuted}\insertinstitute\par
    \vspace{0.15cm}
    \footnotesize\color{textMuted}\insertdate
  \end{beamercolorbox}
\end{frame}

% =============================================================================
% SECCIÓN 1: Problemática Y PERSPECTIVA HISTÓRICA
% =============================================================================
\section{Problemática y Perspectiva Histórica}

% -----------------------------------------------------------------------------
% SLIDE 2: Problemática Central
% -----------------------------------------------------------------------------
\begin{frame}{Problemática Central}{La irrupción de la destreza
  pseudo-cognitiva dual}
  \conceptbox{accentTeal}{La pregunta fundamental}{%
    ¿De qué se compone el intelecto humano cuando los artefactos técnicos
    conjugan simultáneamente el \textbf{diálogo fluido probabilístico} y
    \textbf{cómputo exacto y determinista}?%
  }

  \vspace{0.4cm}
\begin{itemize}\setlength{\itemsep}{0.6em}
    \item La Inteligencia Artificial contemporánea (LLMs) exhibe una
      \textbf{capacidad pseudo-cognitiva dual}:
    \begin{itemize}\setlength{\itemsep}{0.3em}
        \item \textbf{Eficacia}: Síntesis probabilística masiva,
          articulación estilística impecable y automatización de la
          mediación simbólica.
        \item \textbf{Límite}: Carencia total de intencionalidad,
          afectividad, juicio ético encarnado y comprensión fenoménica.
      \end{itemize}
    \item El desafío pedagógico urgente: \textbf{no} prohibir ni
      venerar la herramienta, sino clarificar qué facultades humanas
      delegamos y cuáles debemos fortalecer con mayor rigor.
  \end{itemize}
\end{frame}

% -----------------------------------------------------------------------------
% SLIDE 3: Cronología de la Mediación Externa del Conocimiento
% -----------------------------------------------------------------------------
\begin{frame}{La Fricción Técnica en Perspectiva Histórica}{De la
  pintura rupestre a las redes neuronales}
  \centering
  \resizebox{0.95\textwidth}{!}{
    \begin{tikzpicture}[
        node distance=1.8cm,
        stepnode/.style={rectangle, rounded corners=3pt,
          draw=primaryDark!60, fill=cardBg, line width=0.8pt, inner
        sep=6pt, align=center, font=\footnotesize},
        yearnode/.style={font=\scriptsize\bfseries\color{accentTeal}},
        arrowline/.style={-{Latex[scale=0.8]}, line width=1pt, primaryDark!50}
      ]
      \node[stepnode] (cave) {\textbf{Pinturas Rupestres}\\Memoria
      pictórica externa};
      \node[stepnode, right=0.6cm of cave] (plato)
      {\textbf{Escritura Fonética}\\Advertencia del \textit{Fedro}};
      \node[stepnode, right=0.6cm of plato] (press)
      {\textbf{Imprenta Gutenberg}\\Reproducción mecánica};
      \node[stepnode, right=0.6cm of press] (prussia)
      {\textbf{Escuela Prusiana}\\Estandarización masiva};

      \node[stepnode, below=1.2cm of cave] (calc)
      {\textbf{Calculadora}\\Disputa en matemáticas};
      \node[stepnode, right=0.6cm of calc] (net) {\textbf{Internet /
      Web}\\Memoria hiperdistribuida};
      \node[stepnode, right=0.6cm of net, fill=accentTeal!10,
      draw=accentTeal, line width=1.2pt] (ai) {\textbf{Inteligencia
      Artificial}\\Pseudo-cognición generativa};

      \draw[arrowline] (cave) -- (plato);
      \draw[arrowline] (plato) -- (press);
      \draw[arrowline] (press) -- (prussia);
      \draw[arrowline] (prussia.south) -- ++(0,-0.6) -| (calc.north);
      \draw[arrowline] (calc) -- (net);
      \draw[arrowline] (net) -- (ai);
    \end{tikzpicture}
  }

  \vspace{0.4cm}
  \conceptbox{secondarySlate}{Lección de la Epistemología Histórica}{
    Cada salto técnico genera una crisis pedagógica idéntica: el
    temor a la pérdida de una facultad interna frente a la mediación
    de un dispositivo externo.
  }
\end{frame}

% -----------------------------------------------------------------------------
% SLIDE 4: La Advertencia de Platón en el Fedro
% -----------------------------------------------------------------------------
\begin{frame}{La Memoria Viva frente a la Letra Impresa}{El mito de
  Theuth y Thamus en el \textit{Fedro} de Platón}
  \begin{columns}[T]
    \begin{column}{0.48\textwidth}
      \conceptbox{primaryDark}{El Diagnóstico Platónico}{
        \textit{<<Este invento provocará el olvido en las almas de
          quienes lo aprendan, al descuidar la memoria, ya que
        confiando en lo escrito recordarán desde fuera...>>}
        \par\raggedleft\footnotesize--- Platón, \textit{Fedro} (274e--275b)
      }
    \end{column}
    \begin{column}{0.48\textwidth}
    \begin{itemize}\setlength{\itemsep}{0.6em}
        \item \textbf{El concepto de \textit{Pharmakon}}: La técnica
          es simultáneamente veneno (atrofia de la memoria viva) y
          remedio (preservación archival).
        \item \textbf{\textit{Doxa} vs. \textit{Episteme}}: La
          apariencia de saber (acceso a textos o respuestas
          prefabricadas) versus la apropiación dialéctica del juicio.
        \item \textbf{Paralelismo contemporáneo}: El usuario de LLMs
          que confunde la fluidez de un resumen generado con la
          comprensión real del fenómeno.
      \end{itemize}
    \end{column}
  \end{columns}
\end{frame}

% -----------------------------------------------------------------------------
% SLIDE 5: De Gutenberg al Modelo Prusiano
% -----------------------------------------------------------------------------
\begin{frame}{La Industrialización de la Instrucción}{Imprenta,
  fábrica y escolarización masiva}
\begin{itemize}\setlength{\itemsep}{0.6em}
    \item \textbf{La imprenta de Gutenberg (1440)}: Desacopla la
      replicación del conocimiento del copista manual; democratiza
      el acceso, pero exige alfabetización universal.
    \item \textbf{El modelo industrial prusiano (s. XVIII--XIX)}:
    \begin{itemize}\setlength{\itemsep}{0.3em}
        \item Diseñado para la era industrial: sincronización
          horaria, currículo fragmentado, disciplina y
          estandarización de cohortes.
        \item Premisa subyacente: el ser humano como ejecutor
          confiable de cómputos mecánicos, retención memorística y
          protocolos fijos.
      \end{itemize}
    \item \textbf{La paradoja actual}:
    \begin{itemize}\setlength{\itemsep}{0.3em}
        \item Seguimos evaluando a los estudiantes con métricas
          mecánicas concebidas para la fábrica prusiana.
        \item La IA ejecuta precisamente esas tareas mecánicas de
          manera casi instantánea y a costo marginal cero.
      \end{itemize}
  \end{itemize}
\end{frame}

% -----------------------------------------------------------------------------
% SLIDE 6: Las Controversias Modernas
% -----------------------------------------------------------------------------
\begin{frame}{Calculadoras y Memoria Distribuida}{Precedentes
  directos del debate actual}
  \begin{columns}[T]
    \begin{column}{0.48\textwidth}
      \conceptbox{accentTeal}{La Calculadora Electrónica (1970--1980)}{
        \textbf{El temor}: Atrofia del cálculo mental y la
        aritmética básica.\\[0.4em]
        \textbf{El desenlace}: Liberación de fricción aritmética
        manual para habilitar razonamiento algebraico y modelado de
        sistemas complejos.
      }
    \end{column}
    \begin{column}{0.48\textwidth}
      \conceptbox{accentGold}{La Memoria Web y Google (2000--2010)}{
        \textbf{El temor}: Pérdida de la retención factual (``efecto
        Google'').\\[0.4em]
        \textbf{El desenlace}: Traslado de la competencia: de
        memorizar datos enciclopédicos a evaluar fuentes, indexar y
        sintetizar redes de información.
      }
    \end{column}
  \end{columns}

  \vspace{0.4cm}
  \conceptbox{secondarySlate}{El Salto Cualitativo de la IA Generativa}{
    La calculadora operaba sobre números sintácticos; Internet sobre
    documentos estáticos. La IA procesa y genera \textbf{lenguaje
    natural}, simulando la deliberación y el razonamiento conceptual.
  }
\end{frame}

% -----------------------------------------------------------------------------
% SLIDE 7: La Condición de la IA
% -----------------------------------------------------------------------------
\begin{frame}{La Condición de la Inteligencia Artificial}{Estructura
  y límites de los Modelos de Lenguaje}
  \begin{columns}[T]
    \begin{column}{0.48\textwidth}
      \textbf{\color{primaryDark}Destreza Pseudo-Cognitiva:}
    \begin{itemize}\setlength{\itemsep}{0.4em}
        \item Predicción autoregresiva de tokens condicionada por
          distribución estadística:
          \[
            P(w_t \mid w_1, w_2, \dots, w_{t-1}; \theta)
          \]
        \item Coherencia estilística y fluidez sintáctica sin precedentes.
        \item Capacidad de recombinación intertextual veloz entre
          dominios dispares.
      \end{itemize}
    \end{column}
    \begin{column}{0.48\textwidth}
      \textbf{\color{accentGold}Ausencias Fundamentales:}
    \begin{itemize}\setlength{\itemsep}{0.4em}
        \item \textbf{Cero Anclaje Semántico}: Los símbolos no
          remiten a un mundo fenoménico ni a vivencias somáticas.
        \item \textbf{Cero Intencionalidad}: No hay voluntad,
          compromiso epistémico ni conciencia del error.
        \item \textbf{Inercia hacia la Mediana}: Propensión a
          generar respuestas verosímiles pero conceptualmente
          banales o alucinadas.
      \end{itemize}
    \end{column}
  \end{columns}
\end{frame}

% -----------------------------------------------------------------------------
% SLIDE 8: ¿Qué Compone el Intelecto Humano?
% -----------------------------------------------------------------------------
\begin{frame}{¿Qué Compone el Intelecto Humano?}{Más allá de la
  correlación probabilística}
\begin{itemize}\setlength{\itemsep}{0.6em}
    \item \textbf{Intencionalidad y Sentido}: La capacidad de
      dirigir el pensamiento hacia algo, asumiendo responsabilidad
      por la veracidad de la afirmación.
    \item \textbf{Juicio Crítico y Duda Radical}: La facultad de
      desconfiar de lo autoevidente, cuestionar axiomas establecidos
      y formular contra-hipótesis.
    \item \textbf{Afectividad y Vulnerabilidad}: Comprender las
      consecuencias de una decisión porque se experimenta el
      sufrimiento, el riesgo y el tiempo finito.
    \item \textbf{Imaginación Dialéctica}: Producir saltos
      abductivos que contradicen el patrón histórico previo (ruptura
      de la distribución de entrenamiento).
  \end{itemize}

  \vspace{0.4cm}
  \conceptbox{accentTeal}{Distinción Crucial}{
    Un modelo de lenguaje predice la continuación \textit{más
    probable}. El intelecto innovador busca la verdad o la creación
    \textit{improbable pero fecunda}.
  }
\end{frame}

% -----------------------------------------------------------------------------
% SLIDE 9: Cartografía de la Fricción Epistémica
% -----------------------------------------------------------------------------
\begin{frame}{Cartografía de la Fricción Epistémica}{Fricción
  Mecánica vs. Fricción Reflexiva}
  \begin{columns}[T]
    \begin{column}{0.48\textwidth}
      \conceptbox{accentTeal}{Fricción Mecánica (Delegable)}{
      \begin{itemize}\setlength{\itemsep}{0.3em}
          \item Búsqueda rutinaria de referencias bibliográficas.
          \item Reescritura sintáctica y traducción de formatos.
          \item Depuración de errores tipográficos y sintaxis de código.
          \item Resumen inicial de extensos corpus documentales.
        \end{itemize}
        \vspace{0.2em}
        \textit{Resultado}: Ahorro de tiempo sin pérdida de
        profundidad conceptual.
      }
    \end{column}
    \begin{column}{0.48\textwidth}
      \conceptbox{primaryDark}{Fricción Reflexiva (Intransferible)}{
      \begin{itemize}\setlength{\itemsep}{0.3em}
          \item Formulación de la pregunta de investigación genuina.
          \item Detección de sesgos metodológicos y éticos.
          \item Resolución de paradojas y contradicciones teóricas.
          \item Toma de decisiones bajo incertidumbre radical.
        \end{itemize}
        \vspace{0.2em}
        \textit{Resultado}: La zona donde ocurre la verdadera
        maduración intelectual.
      }
    \end{column}
  \end{columns}
\end{frame}

% -----------------------------------------------------------------------------
% SLIDE 10: Neutralidad de Doctrinas Pedagógicas
% -----------------------------------------------------------------------------
\begin{frame}{Neutralidad de Doctrinas Pedagógicas}{Superando falsas
  dicotomías metodológicas}
  \conceptbox{secondarySlate}{Postura de la Conferencia}{
    No se defiende un dogma pedagógico excluyente (conductismo vs.
    constructivismo vs. conectivismo), sino un \textbf{análisis
    epistémico riguroso} de las condiciones del aprendizaje.
  }

  \vspace{0.4cm}
\begin{itemize}\setlength{\itemsep}{0.6em}
    \item \textbf{Contra el rechazo tecnofóbico}: Ignorar la IA
      produce una educación anacrónica y desconectada de los medios
      reales de producción de conocimiento.
    \item \textbf{Contra el optimismo acrítico}: Someter la
      pedagogía a interfaces algorítmicas comerciales conduce a la
      atrofia del pensamiento independiente.
    \item \textbf{El criterio de arbitraje}: La pedagogía debe
      medirse por el grado de \textbf{agencia, discernimiento y
      capacidad innovadora} que despierta en el aprendiz.
  \end{itemize}
\end{frame}

% -----------------------------------------------------------------------------
% SLIDE 11: La IA como Arnés Metodológico
% -----------------------------------------------------------------------------
\begin{frame}{La IA como Arnés Metodológico}{Soporte de tracción sin
  sustitución de agencia}
  \begin{columns}[T]
    \begin{column}{0.5\textwidth}
      \conceptbox{accentTeal}{¿Por qué un ``Arnés''?}{
        Un arnés no camina por el montañista; le confiere seguridad,
        absorbe la tensión mecánica del impacto y le permite
        alcanzar cumbres que de otro modo serían inaccesibles.
      }
      \vspace{0.3cm}
    \begin{itemize}\setlength{\itemsep}{0.3em}
        \item \textbf{Exploración rápida}: Mapear paisajes
          conceptuales en minutos.
        \item \textbf{Contra-argumentador}: Usar el modelo como
          interlocutor socrático imperfecto.
      \end{itemize}
    \end{column}
    \begin{column}{0.46\textwidth}
    \begin{itemize}\setlength{\itemsep}{0.5em}
        \item \textbf{Retroalimentación inmediata}: Iteración
          temprana sobre borradores e hipótesis de trabajo.
        \item \textbf{Exigencia al estudiante}: El estudiante debe
          justificar por qué acepta o rechaza cada output emitido
          por la máquina.
        \item \textbf{Desacoplamiento}: El valor no reside en la
          respuesta rápida del LLM, sino en la \textit{crítica
          sistemática} de esa respuesta.
      \end{itemize}
    \end{column}
  \end{columns}
\end{frame}

% -----------------------------------------------------------------------------
% SLIDE 12: Descarga Cognitiva (Cognitive Offloading)
% -----------------------------------------------------------------------------
\begin{frame}{Descarga Cognitiva (\textit{Cognitive Offloading})}{El
  filo de la navaja: entre la liberación y la atrofia}
  \begin{columns}[T]
    \begin{column}{0.48\textwidth}
      \conceptbox{accentTeal}{El Beneficio Prometido}{
        Descargar operaciones de memoria a corto plazo y cálculo
        procedimental para reinvertir el ancho de banda biológico en
        tareas meta-cognitivas de orden superior.
      }
    \end{column}
    \begin{column}{0.48\textwidth}
      \conceptbox{accentGold}{El Riesgo Sistémico}{
        Cuando se descarga el esfuerzo de interpretación, el cerebro
        no reasigna automáticamente los recursos a la creatividad:
        puede caer en la \textbf{inercia cognitiva} y la complacencia pasiva.
      }
    \end{column}
  \end{columns}

  \vspace{0.4cm}
\begin{itemize}\setlength{\itemsep}{0.4em}
    \item \textbf{La paradoja de la facilidad}: El aprendizaje
      profundo requiere cierto nivel de dificultad deseable
      (\textit{desirable difficulties}). Si se elimina toda
      fricción, el aprendizaje se vuelve superficial y volátil.
  \end{itemize}
\end{frame}

% -----------------------------------------------------------------------------
% SLIDE 16: Diagrama Conceptual TikZ: El Ciclo Pedagógico con Arnés
% -----------------------------------------------------------------------------
\begin{frame}{El Ciclo Pedagógico con Arnés}{Reconfiguración del
  proceso de aprendizaje}
  \centering
  \begin{tikzpicture}[
      scale=0.9, every node/.style={transform shape},
      boxnode/.style={rectangle, rounded corners=4pt,
        draw=primaryDark, fill=cardBg, line width=1pt, inner sep=8pt,
      align=center, text width=3.0cm},
      arnesnode/.style={rectangle, rounded corners=4pt,
        draw=accentGold, fill=accentGold!10, line width=1.2pt, inner
      sep=8pt, align=center, text width=3.2cm},
      conatusnode/.style={rectangle, rounded corners=4pt,
        draw=accentTeal, fill=accentTeal!12, line width=1.4pt, inner
      sep=8pt, align=center, text width=3.4cm},
      flowarrow/.style={-{Latex[scale=0.9]}, line width=1.2pt, primaryDark!70}
    ]
    \node[boxnode] (q) at (0, 2.2) {\textbf{1.
    Formulación}\\Pregunta epistémica y contextualización};
    \node[arnesnode] (explore) at (4.5, 2.2) {\textbf{2. Arnés
    IA}\\Fricción mecánica:\\resumen, código, borrador};
    \node[boxnode] (audit) at (4.5, -1.2) {\textbf{3. Auditoría
    Crítica}\\Detección de sesgos, falsedades y límites};
    \node[conatusnode] (conatus) at (0, -1.2) {\textbf{4.
    \textit{Conatus ad Novandum}}\\Innovación reflexiva y síntesis humana};

    \draw[flowarrow] (q) -- node[above,
    font=\scriptsize\color{textMuted}] {disparo} (explore);
    \draw[flowarrow] (explore) -- node[right,
    font=\scriptsize\color{textMuted}] {evaluación} (audit);
    \draw[flowarrow] (audit) -- node[below,
    font=\scriptsize\color{textMuted}] {subsunción} (conatus);
    \draw[flowarrow] (conatus) -- node[left,
    font=\scriptsize\color{textMuted}] {nueva pregunta} (q);
  \end{tikzpicture}

  \vspace{0.2cm}
  \footnotesize\color{textMuted}
  \textit{El ciclo cierra en la re-problematización continua: la
    máquina amortigua el paso mecánico (2), pero la apertura (1), la
  duda (3) y la creación (4) corresponden al agente humano.}
\end{frame}

% -----------------------------------------------------------------------------
% SLIDE 14: Reasignación Consciente de Recursos Cognitivos
% -----------------------------------------------------------------------------
\begin{frame}{Reasignación Consciente de Recursos}{De la ejecución
  repetitiva a la contemplación activa}
\begin{itemize}\setlength{\itemsep}{0.6em}
    \item \textbf{Agotamiento procedimental tradicional}:
      Históricamente, hasta el $80\%$ del esfuerzo estudiantil se
      consumía en tareas instrumentales (formateo, búsqueda
      bibliográfica elemental, resolución algebraica estándar).
    \item \textbf{La oportunidad de la reasignación}:
    \begin{itemize}\setlength{\itemsep}{0.3em}
        \item \textbf{Validez metodológica}: ¿por qué esta premisa y no otra?
        \item \textbf{Visión transdisciplinaria}: conectar física, ética,
          biología y computación.
        \item \textbf{Imaginación contrafáctica}: ¿qué sucedería si el
          supuesto fundamental fuera falso?
      \end{itemize}
    \item \textbf{La condición indispensable}: La reasignación no
      ocurre por defecto; requiere un diseño curricular explícito
      que premie la originalidad y penalice la mera agregación sintáctica.
  \end{itemize}
\end{frame}

% -----------------------------------------------------------------------------
% SLIDE 15: Concepto Filosófico: Conatus ad Novandum
% -----------------------------------------------------------------------------
\begin{frame}{El Concepto de \textit{Conatus ad Novandum}}{El núcleo
  del impulso humano}
  \begin{columns}[T]
    \begin{column}{0.48\textwidth}
      \conceptbox{accentTeal}{La Raíz Spinoziana (\textit{Conatus})}{
        En Spinoza, el \textit{conatus} es la fuerza o conato
        inherente a cada ser para perseverar en su existencia y
        potenciar su capacidad de obrar (\textit{potentia agendi}).
      }
    \end{column}
    \begin{column}{0.48\textwidth}
      \conceptbox{primaryDark}{El Giro: \textit{ad Novandum}}{
        No se trata únicamente de perseverar en el estado presente,
        sino del \textbf{impulso deliberado por renovar, crear,
        desestabilizar lo dado y fundar nuevas perspectivas}.
      }
    \end{column}
  \end{columns}

  \vspace{0.4cm}
\begin{itemize}\setlength{\itemsep}{0.4em}
    \item La máquina persevera en su distribución estadística
      (inercia a la media del corpus).
    \item El ser humano manifiesta su vitalidad cuando se atreve a
      quebrar la regla para alumbrar un sentido inédito:
      \textbf{\textit{Conatus ad Novandum}}.
  \end{itemize}
\end{frame}

% -----------------------------------------------------------------------------
% SLIDE 16: Formación de Líderes Visionarios
% -----------------------------------------------------------------------------
\begin{frame}{Formación de Líderes Visionarios}{De operadores de
  respuestas a arquitectos de preguntas}
  \begin{columns}[T]
    \begin{column}{0.48\textwidth}
      \textbf{\color{textMuted}El Operador de Algoritmos:}
    \begin{itemize}\setlength{\itemsep}{0.3em}
        \item Acepta las opciones presentadas por el sistema.
        \item Optimiza métricas predefinidas de eficiencia.
        \item Delega el criterio de verdad en la máquina.
        \item Reproduce los sesgos del pasado con mayor celeridad.
      \end{itemize}
    \end{column}
    \begin{column}{0.48\textwidth}
      \textbf{\color{primaryDark}El Líder Visionario Humano:}
    \begin{itemize}\setlength{\itemsep}{0.3em}
        \item Define \textit{qué} problemas merecen ser resueltos y por qué.
        \item Juzga el impacto ético y social sobre personas concretas.
        \item Tolera la ambigüedad y la contradicción fecunda.
        \item Inspira a través del testimonio, la empatía y la convicción.
      \end{itemize}
    \end{column}
  \end{columns}

  \vspace{0.4cm}
  \conceptbox{secondarySlate}{Máxima para el Liderazgo Educativo}{
    Si la IA abarata las respuestas, el valor supremo reside en la
    formulación de las \textbf{preguntas decisivas}.
  }
\end{frame}

% -----------------------------------------------------------------------------
% SLIDE 17: Transformación del Rol Docente
% -----------------------------------------------------------------------------
\begin{frame}{Transformación del Rol Docente}{De transmisor de
  información a maestro del juicio}
\begin{itemize}\setlength{\itemsep}{0.6em}
    \item \textbf{La muerte definitiva de la cátedra magistral
      transmisiva}: El docente como mero dispensador enciclopédico
      de datos ha sido superado por la indexación algorítmica.
    \item \textbf{La emergencia del tutor dialéctico y socrático}:
    \begin{itemize}\setlength{\itemsep}{0.3em}
        \item Desafiar las respuestas generadas por los estudiantes
          y sus herramientas de IA.
        \item Modelar el arte del escepticismo metodológico y la
          honestidad intelectual.
        \item Guiar la atención en un ecosistema saturado de
          estímulos hiper-estimulantes.
      \end{itemize}
    \item \textbf{Evaluación basada en el proceso}: No evaluar el
      artefacto final estático (el ensayo o el código que un LLM
      genera en segundos), sino la \textbf{cadena de decisiones,
      correcciones y argumentaciones} sostenidas por el estudiante.
  \end{itemize}
\end{frame}

% -----------------------------------------------------------------------------
% SLIDE 18: El Dilema Epistemológico: A priori vs. A posteriori
% -----------------------------------------------------------------------------
\begin{frame}{El Dilema Epistemológico}{Estructura formal abstracta
  (\textit{a priori}) vs. Afectación somática (\textit{a posteriori})}
  \begin{columns}[T]
    \begin{column}{0.48\textwidth}
      \conceptbox{primaryDark}{El Polo \textit{A Priori} / Formal}{
      \begin{itemize}\setlength{\itemsep}{0.3em}
          \item Estructuras matemáticas, relaciones lógicas puras y
            pesos relacionales en grafos de alta dimensión.
          \item Cómputo incorpóreo: invariante frente a la biología
            o la temperatura emocional.
          \item Donde la IA opera con soberbia velocidad combinatoria.
        \end{itemize}
      }
    \end{column}
    \begin{column}{0.48\textwidth}
      \conceptbox{accentTeal}{El Polo \textit{A Posteriori} / Somático}{
      \begin{itemize}\setlength{\itemsep}{0.3em}
          \item La experiencia fenoménica del mundo a través de
            receptores sensoriales y dolor/placer somático.
          \item El aprendizaje biológico condicionado por la finitud
            y la vulnerabilidad física.
          \item La fuente primaria del significado existencial humano.
        \end{itemize}
      }
    \end{column}
  \end{columns}

  \vspace{0.4cm}
  \conceptbox{secondarySlate}{El Vacío de la Inteligencia Incorpórea}{
    Un sistema que solo manipula relaciones sintácticas carece de
    acceso directo al mundo fenoménico: habita exclusivamente el
    reino de la correlación formal desincorporada.
  }
\end{frame}

% -----------------------------------------------------------------------------
% SLIDE 19: La Dialéctica: Cogito vs. Sentio
% -----------------------------------------------------------------------------
\begin{frame}{La Dialéctica: \textit{Cogito} vs. \textit{Sentio}}{La
  mente abstracta frente a la corporalidad encarnada}
  \centering
  \begin{tikzpicture}[
      scale=0.9, every node/.style={transform shape},
      boxA/.style={rectangle, rounded corners=5pt, draw=primaryDark,
        fill=primaryDark!8, line width=1.2pt, inner sep=10pt, text
      width=4.5cm, align=center},
      boxB/.style={rectangle, rounded corners=5pt, draw=accentGold,
        fill=accentGold!10, line width=1.2pt, inner sep=10pt, text
      width=4.5cm, align=center},
      bridge/.style={rectangle, rounded corners=3pt,
        draw=accentTeal, fill=cardBg, line width=1pt, inner sep=6pt,
      text width=3.0cm, align=center, font=\footnotesize}
    ]
    \node[boxA] (cogito) at (0, 0) {
      \Large\textbf{\color{primaryDark}\textit{COGITO}}\\[0.4em]
      \small
      Manipulación simbólica\\
      Abstracción algorítmica\\
      Cálculo formal incorpóreo\\
      Optimización de funciones
    };

    \node[boxB] (sentio) at (8, 0) {
      \Large\textbf{\color{accentGold}\textit{SENTIO}}\\[0.4em]
      \small
      Experiencia sensorial directa\\
      Afectividad y vulnerabilidad\\
      Cognición encarnada (\textit{embodied})\\
      Conciencia fenomenológica
    };

    \draw[<->, line width=1.5pt, primaryDark!60] (cogito) --
    node[above=0.2cm, bridge] {Tensión Irreductible\\de la
    Existencia} (sentio);
  \end{tikzpicture}

  \vspace{0.4cm}
  \conceptbox{secondarySlate}{Implicación Pedagógica}{
    La educación tradicional privilegia el \textit{cogito} formal
    desvinculado; la IA nos fuerza a recordar que la inteligencia
    genuina está irreductiblemente enraizada en el \textit{sentio} vivo.
  }
\end{frame}

% -----------------------------------------------------------------------------
% SLIDE 20: ¿Puede existir un ``Animal Cognitivo'' Incorpóreo?
% -----------------------------------------------------------------------------
\begin{frame}{El Mito del Intelecto Desencarnado}{¿Existe un
  ``animal cognitivo'' o es solo un artefacto matemático?}
\begin{itemize}\setlength{\itemsep}{0.6em}
    \item \textbf{La falacia cartesiana reeditada}: La idea de que
      la mente puede desacoplarse del organismo biológico sin perder
      su esencia constitutiva.
    \item \textbf{La perspectiva de la cognición encarnada
      (\textit{Embodied Cognition})}:
    \begin{itemize}\setlength{\itemsep}{0.3em}
        \item No pensamos \textit{a pesar} de tener cuerpo, sino
          \textit{a través} del cuerpo: los conceptos abstractos son
          metáforas somatosensoriales (Varela, Thompson, Lakoff).
        \item La cognición biológica emergió para regular la acción
          y la supervivencia en un entorno termodinámico adverso, no
          para resolver acertijos lógicos aislados.
      \end{itemize}
    \item \textbf{El veredicto sobre la IA pura}:
    \begin{itemize}\setlength{\itemsep}{0.3em}
        \item La máquina no es un ``animal cognitivo''; es un
          \textbf{espejo estadístico} de los registros dejados por
          seres que sí poseían cuerpos y vivencias.
      \end{itemize}
  \end{itemize}
\end{frame}

% -----------------------------------------------------------------------------
% SLIDE 21: Abstracciones Simbólicas y Manipulación Ambiental
% -----------------------------------------------------------------------------
\begin{frame}{Abstracciones Simbólicas y Manipulación Ambiental}{La
  instrumentalización al servicio de la conveniencia y la codicia}
  \begin{columns}[T]
    \begin{column}{0.48\textwidth}
      \conceptbox{primaryDark}{El Origen de la Abstracción}{
        La abstracción simbólica ha sido la herramienta adaptativa
        suprema del \textit{Homo sapiens}: permitió codificar
        recursos, anticipar estaciones y coordinar tribus a gran escala.
      }
    \end{column}
    \begin{column}{0.48\textwidth}
      \conceptbox{accentGold}{La Deriva Instrumental}{
        Al independizarse del límite ecológico y somático, la
        abstracción se convirtió en motor de dominación y extracción
        desmedida en beneficio de la codicia y la conveniencia inmediata.
      }
    \end{column}
  \end{columns}

  \vspace{0.4cm}
\begin{itemize}\setlength{\itemsep}{0.4em}
    \item La IA representa la culminación técnica de esta
      abstracción instrumental: procesar datos del mundo real a
      espaldas de la fragilidad del ecosistema y de los sujetos.
    \item La educación no puede ser mero entrenamiento para servir a
      esta maquinaria de conveniencia.
  \end{itemize}
\end{frame}

% -----------------------------------------------------------------------------
% SLIDE 22: Evolución No-Teleológica y Bifurcación
% -----------------------------------------------------------------------------
\begin{frame}{Evolución No-Teleológica y Bifurcación}{¿Hacia una
  escisión entre la biología y la técnica algorítmica?}
\begin{itemize}\setlength{\itemsep}{0.6em}
    \item \textbf{Evolución sin diseño previo}:
    \begin{itemize}\setlength{\itemsep}{0.3em}
        \item La evolución biológica no posee meta prefijada
          (\textit{telos}); opera por deriva genética, presiones
          ambientales y contingencia histórica.
        \item La cultura y la técnica son extensiones epigenéticas
          que retroalimentan las presiones de selección.
      \end{itemize}
    \item \textbf{La hipótesis de la bifurcación evolutiva}:
    \begin{itemize}\setlength{\itemsep}{0.3em}
        \item ¿Estamos presenciando una bifurcación entre la estirpe
          somática biológica y un ecosistema exógeno de agentes
          algorítmicos auto-perfeccionados?
        \item O bien, ¿una absorción en la cual el ser humano adapta
          su cognición a los formatos digestibles por los algoritmos
          (empobrecimiento lingüístico y cognitivo)?
      \end{itemize}
    \item \textbf{El papel de la educación}: Actuar como salvaguarda
      evolutiva que preserve la riqueza fenoménica y la autonomía
      crítica del organismo humano.
  \end{itemize}
\end{frame}

% -----------------------------------------------------------------------------
% SLIDE 23: La Pregunta Límite: ¿Qué es la Vida?
% -----------------------------------------------------------------------------
\begin{frame}{La Pregunta Límite: ¿Qué es la
  \textit{Vida}?}{Frontera ontológica entre el ser autopoiético y el
  artefacto técnico}
  \begin{columns}[T]
    \begin{column}{0.48\textwidth}
      \conceptbox{accentTeal}{El Organismo Vivo (\textit{Autopoiesis})}{
      \begin{itemize}\setlength{\itemsep}{0.3em}
          \item Regeneración celular continua contra el gradiente entrópico.
          \item Sensibilidad al dolor, a la muerte y al tiempo irreversible.
          \item Significado como necesidad vital y supervivencia situada.
        \end{itemize}
      }
    \end{column}
    \begin{column}{0.48\textwidth}
      \conceptbox{primaryDark}{El Artefacto Sintético (Algoritmo)}{
      \begin{itemize}\setlength{\itemsep}{0.3em}
          \item Reversible en sus operaciones de memoria y parámetros.
          \item Indiferente a su apagado o reconfiguración física.
          \item Sintaxis sin semántica, computación sin dolor ni regocijo.
        \end{itemize}
      }
    \end{column}
  \end{columns}

  \vspace{0.4cm}
  \conceptbox{secondarySlate}{El Destino del Pensador}{
    Cultivar futuros pensadores significa educar seres que
    reconozcan el valor irreemplazable de la \textbf{vida viva}
    frente a cualquier simulación sintética deslumbrante.
  }
\end{frame}

% -----------------------------------------------------------------------------
% SLIDE 24: Matriz Comparativa de Tensiones Epistémicas
% -----------------------------------------------------------------------------
\begin{frame}{Matriz de Tensiones Epistémicas}{Síntesis del marco conceptual}
  \centering
  \scriptsize
  \renewcommand{\arraystretch}{1.3}
  \begin{tabular}{p{3.0cm} p{4.5cm} p{5.0cm}}
    \toprule
    \textbf{Dimensión} & \textbf{Lógica Algorítmica (IA)} &
    \textbf{Formación Humana (\textit{Conatus})} \\
    \midrule
    \textbf{Operación Nuclear} & Optimización probabilística y
    predicción de la media & Desvío crítico, asombro y salto
    abductivo (\textit{novandum}) \\
    \textbf{Fricción Tratada} & Absorción y supresión de la fricción
    mecánica & Cultivo pedagógico de la fricción reflexiva \\
    \textbf{Eje Epistémico} & \textit{Cogito} abstracto
    desincorporado y simbólico & Dialéctica encarnada:
    \textit{Cogito} arraigado en el \textit{Sentio} \\
    \textbf{Relación con el Saber} & Agregación intertextual y
    paráfrasis probabilística & Juicio fundado en la experiencia, la
    ética y la finitud \\
    \textbf{Horizonte de Acción} & Eficiencia instrumental y
    automatización & Liderazgo con sentido, justicia y creatividad genuina \\
    \bottomrule
  \end{tabular}
\end{frame}

% -----------------------------------------------------------------------------
% SLIDE 25: Preguntas Abiertas para el Diálogo
% -----------------------------------------------------------------------------
\begin{frame}{Preguntas Abiertas para el Diálogo}{¿Hacia dónde
  cultivamos a los futuros pensadores?}
\begin{enumerate}\setlength{\itemsep}{0.8em}
    \item Si la máquina asume la articulación del discurso formal,
      \textbf{¿cómo rediseñamos las aulas para que los estudiantes
        no pierdan la capacidad de articular pensamientos propios sin
      mediación algorítmica?}
    \item ¿Podemos construir un sistema educativo que aproveche el
      \textbf{arnés técnico} de la IA para liberar tiempo sin que
      ello derive en la atrofia del esfuerzo reflexivo?
    \item Frente a la tentación de reducir la educación a
      competencias instrumentales para el mercado tecnológico,
      \textbf{¿cómo recuperamos el \textit{Conatus ad Novandum} como
      fin en sí mismo?}
    \item En la era de la inteligencia incorpórea: \textbf{¿qué
      significa salvaguardar lo humano?}
  \end{enumerate}

  \vspace{0.4cm}
  \centering
  \small\textbf{\color{primaryDark}¡Muchas gracias por su atención y diálogo!}
\end{frame}

\end{document}
